
\subsubsection{Interpreting the Detection Signal}

The $8\times$ loss difference provides a statistically robust signal for minority detection. However, translating distributional difference into individual identification is constrained by base rate mathematics: at 5\% minority rate, if we predict 5\% of samples as minority (to match the true rate), perfect ranking achieves approximately 33\% precision. The achieved 32.9\% precision approaches this bound.

This does not indicate weak signal---the Mann-Whitney p-value of 7.36 $\times 10^{-15}$ demonstrates extremely strong distributional separation. Rather, it reflects the fundamental challenge of rare class detection. The appropriate use case is weighted training (as in JTT) rather than hard minority labeling.

\subsubsection{Magnitude Gap versus Timing Gap}

The original hypothesis predicted that spurious features would peak before core features during training, reflecting their ''easier'' status. With ImageNet-pretrained features, this prediction was not supported---both feature types peaked at epoch 81. Simplicity bias manifests as consistently higher spurious encoding (9-11\% gap) rather than earlier encoding.

This finding may be specific to transfer learning conditions. ImageNet pretraining likely encodes both background and object patterns in the initial representations, compressing the learning dynamics. Training from random initialization may exhibit different temporal patterns, though this was not tested.

\subsubsection{Limitations}

\textbf{Single dataset}: Only Waterbirds was tested. Transfer to CelebA or MultiNLI was planned but not executed.

\textbf{Pretrained features only}: The timing dynamics may differ when training from scratch.

\textbf{No intervention evaluation}: This work characterizes the detection signal but does not evaluate whether upweighting detected samples improves worst-group accuracy. JTT's demonstrated effectiveness with misclassification-based detection suggests promise, but this remains unverified for loss-magnitude detection.

\textbf{Base rate sensitivity}: The precision bounds depend strongly on minority prevalence. At higher minority rates (e.g., 15\% in CelebA), higher absolute precision would be achievable.

